\section{Monoidal categorification of cluster algebras}
\label{Sec : monoidal categorification of cluster algebras}

\subsection{Cluster algebras}
\label{ssec : cluster algebras}

Let $J = J_{\mathrm{uf}} \sqcup J_{\mathrm{fr}}$ be a finite set.
An integer matrix
\[
\widetilde{B}
= \begin{pmatrix} B \\ P \end{pmatrix}
= (b_{ij})_{(i, j) \in J \times J_{\mathrm{uf}}}
\]
is called an \emph{exchange matrix}
if its principal part $B$ is skew-symmetric.
Let $\Lambda = (\lambda_{ij})_{i, j \in J}$
be a skew-symmetric integer matrix.
Let $D = \operatorname{diag}(d_i)_{i \in J_{\mathrm{uf}}}$
be a diagonal integer matrix with $d_i > 0$.
The pair $(\widetilde{B}, \Lambda)$
is called a \emph{compatible pair of type $D$} if
\[
\widetilde{B}^{\mathsf{T}}\Lambda = (D\ 0),
\qquad
(D\ 0) = (\delta_{ij}d_i)_{(i, j) \in J_{\mathrm{uf}} \times J}.
\]
For a compatible pair $(\widetilde{B}, \Lambda)$ of type $D$
and $k \in J_{\mathrm{uf}}$, we define
\[
\mu_k\widetilde{B}
= (b'_{ij})_{(i, j) \in J \times J_{\mathrm{uf}}},
\qquad
\mu_k\Lambda = (\lambda'_{ij})_{i, j \in J}
\]
by
\[
b'_{ij} =
\begin{cases}
    -b_{ij}, & \text{if $i = k$ or $j = k$}, \\
    b_{ij} + [b_{ik}]_+[b_{kj}]_+
    - [-b_{ik}]_+[-b_{kj}]_+, & \text{otherwise},
\end{cases}
\]
\[
\lambda'_{ij} =
\begin{cases}
    -\lambda_{ik}
    + \displaystyle\sum_{\ell \in J}[-b_{\ell k}]_+\lambda_{i\ell},
    & \text{if $j = k$, $i \neq k$}, \\
    -\lambda_{kj}
    + \displaystyle\sum_{\ell \in J}[-b_{\ell k}]_+\lambda_{\ell j},
    & \text{if $i = k$, $j \neq k$}, \\
    \lambda_{ij}, & \text{otherwise},
\end{cases}
\]
where $[a]_+ \coloneq \max\{a, 0\}$ for any $a \in \mathbb{R}$.
It is known that
\begin{itemize}
    \item
    $\mu_k(\widetilde{B}, \Lambda)
    \coloneq (\mu_k\widetilde{B}, \mu_k\Lambda)$
    is also a compatible pair of type $D$,
    i.e., the principal part of $\mu_k\widetilde{B}$ is skew-symmetric,
    $\mu_k\Lambda$ is skew-symmetric, and
    $(\mu_k\widetilde{B})^{\mathsf{T}}\mu_k\Lambda = (D\ 0)$,
    \item $\mu_k^2 = \id$, 
\end{itemize}
see \cite{BZ05}.

Let $\mathbb{F}$ be the rational function field
in $|J|$ variables over $\Q$.
A \emph{seed} is a pair $(\mathbf{x}, \widetilde{B})$ where
\begin{itemize}
    \item $\mathbf{x} = (x_i)_{i \in J}$ satisfies that
    $\{x_i \mid i \in J\}$ is a free generating set
    of $\mathbb{F}$ over $\Q$,
    \item $\widetilde{B}$ is an exchange matrix.
\end{itemize}
For a seed $(\mathbf{x}, \widetilde{B})$ and $k \in J_{\mathrm{uf}}$,
the \emph{mutation} of $(\mathbf{x}, \widetilde{B})$ at $k$
is the pair $(\mu_k\mathbf{x}, \mu_k\widetilde{B})$ given by
\[
(\mu_k\mathbf{x})_i =
\begin{cases}
    x_i, & \text{if $i \neq k$}, \\
    x_k^{-1}\left(
    \displaystyle\prod_{j \in J}x_j^{[b_{jk}]_+}
    + \displaystyle\prod_{j \in J}x_j^{[-b_{jk}]_+}
    \right), & \text{if $i = k$}.
\end{cases}
\]
Then
$\mu_k(\mathbf{x}, \widetilde{B})
\coloneq (\mu_k\mathbf{x}, \mu_k\widetilde{B})$
is also a seed and $\mu_k^2 = \id$.
Let $\mathbb{T}$ denote the $|J_{\mathrm{uf}}|$-regular tree.
We label the edges of $\mathbb{T}$ by $J_{\mathrm{uf}}$
such that the $|J_{\mathrm{uf}}|$ different edges adjacent to the same vertex
of $\mathbb{T}$ receive different labels.

\begin{Def}
\label{Def : cluster patterns}
A \emph{cluster pattern}
$\mathcal{S}_{\mathbf{x}}
= \{(\mathbf{x}_t, \widetilde{B}_t) \mid t \in \mathbb{T}\}$
is an assignment of a seed $(\mathbf{x}_t, \widetilde{B}_t)$
to every vertex $t \in \mathbb{T}$ such that
\[
(\mathbf{x}_{t'}, \widetilde{B}_{t'})
= \mu_k(\mathbf{x}_t, \widetilde{B}_t)
\]
whenever $\begin{xy}(0,1)*+{t}="A",(10,1)*+{t'}="B",\ar@{-}^k"A";"B" \end{xy}$ in $\mathbb{T}$.
A \emph{$\Lambda$-cluster pattern}
$\mathcal{S}_{\mathbf{x}}^{\Lambda}
= \{(\mathbf{x}_t, \widetilde{B}_t, \Lambda_t)
\mid t \in \mathbb{T}\}$ of type $D$
is an assignment of a triple $(\mathbf{x}_t, \widetilde{B}_t, \Lambda_t)$
to every vertex $t \in \mathbb{T}$ such that
$\{(\mathbf{x}_t, \widetilde{B}_t) \mid t \in \mathbb{T}\}$
is a cluster pattern,
$(\widetilde{B}_t, \Lambda_t)$ is a compatible pair of type $D$, and
\[
(\mathbf{x}_{t'}, \widetilde{B}_{t'}, \Lambda_{t'})
= (\mu_k\mathbf{x}_t, \mu_k\widetilde{B}_t, \mu_k\Lambda_t)
\]
whenever $\begin{xy}(0,1)*+{t}="A",(10,1)*+{t'}="B",\ar@{-}^k"A";"B" \end{xy}$ in $\mathbb{T}$.
We call $(\mathbf{x}_t, \widetilde{B}_t, \Lambda_t)$
a \emph{$\Lambda$-seed}.
\end{Def}

For a $\Lambda$-seed $(\mathbf{x}_t, \widetilde{B}_t, \Lambda_t)$,
we usually write
\[
\mathbf{x}_t = (x_{i;t})_{i \in J},
\qquad
\widetilde{B}_t = (b_{ij}^{t}),
\qquad
\Lambda_t = (\lambda_{ij}^{t}).
\]
We call each $x_{i;t}$ $(i \in J)$ a \emph{cluster variable}.
More precisely, we call $x_{i;t}$ $(i \in J_{\mathrm{uf}})$
an \emph{unfrozen cluster variable} and
$x_{i;t}$ $(i \in J_{\mathrm{fr}})$ a \emph{frozen cluster variable}.
A monomial in $\{x_{i;t} \mid i \in J\}$
is called a \emph{cluster monomial}.

\begin{Def}
\label{Def : cluster algebras}
Let
$\mathcal{S}_{\mathbf{x}}^{\Lambda}
= \{(\mathbf{x}_t, \widetilde{B}_t, \Lambda_t)
\mid t \in \mathbb{T}\}$
be a $\Lambda$-cluster pattern.
The \emph{(partially compactified) cluster algebra}
$\overline{\mathcal{A}}(\mathcal{S}_{\mathbf{x}})$
is the $\Z$-subalgebra of $\mathbb{F}$
generated by $\{x_{i;t} \mid i \in J,\ t \in \mathbb{T}\}$.
The \emph{upper cluster algebra}
$\mathcal{U}(\mathcal{S}_{\mathbf{x}})$
is the $\Z$-subalgebra of $\mathbb{F}$ such that
\[
\mathcal{U}(\mathcal{S}_{\mathbf{x}})
= \bigcap_{t \in \mathbb{T}}
\Z[x_{i;t}^{\pm 1} \mid i \in J].
\]
The \emph{(partially compactified) $\Lambda$-cluster algebra}
$\overline{\mathcal{A}}(\mathcal{S}_{\mathbf{x}}^{\Lambda})$
is a cluster algebra
$\overline{\mathcal{A}}(\mathcal{S}_{\mathbf{x}})$
endowed with the additional data
$\Lambda = \{\Lambda_t \mid t \in \mathbb{T}\}$.
The \emph{$\Lambda$-upper cluster algebra}
$\mathcal{U}(\mathcal{S}_{\mathbf{x}}^{\Lambda})$
is defined in the same manner.
\end{Def}

It is known that
$\overline{\mathcal{A}}(\mathcal{S}_{\mathbf{x}})
\subset \mathcal{U}(\mathcal{S}_{\mathbf{x}})$ \cite{FZ02}.
For an exchange matrix $\widetilde{B}$,
we can associate a quiver $Q_{\widetilde{B}}$
characterized by the following properties:
\begin{enumerate}
    \item[(1)] its vertices are labeled by $J$,
    and its frozen vertices are labeled by $J_{\mathrm{fr}}$,
    \item[(2)] it has no loops, $2$-cycles,
    and arrows between frozen vertices,
    \item[(3)] for any $i \in J$ and $j \in J_{\mathrm{uf}}$,
    \[
    b_{ij}
    = (\text{the number of arrows from $i$ to $j$})
    - (\text{the number of arrows from $j$ to $i$}).
    \]
\end{enumerate}
Conversely, every quiver satisfying conditions (1) and (2)
determines an integer matrix $\widetilde{B}$ by condition (3). 

\subsection{Pointed elements}
\label{ssec : pointed elements}

Let
$\mathcal{S}_{\mathbf{x}}^{\Lambda}
= \{(\mathbf{x}_t, \widetilde{B}_t, \Lambda_t)
\mid t \in \mathbb{T}\}$
be a $\Lambda$-cluster pattern.
For $t \in \mathbb{T}$ and $i \in J_{\mathrm{uf}}$, we set
\[
\widehat{y}_{i;t} \coloneq \mathbf{x}_t^{\widetilde{B}_t e_i}.
\]

\begin{Def}[{\cite[Definition 3.1.4]{Qin17},
\cite[Definition 3.2.2]{Qin24}}]
\label{Def : pointed and bipointed elements}
Let $u \in \mathcal{U}(\mathcal{S}_{\mathbf{x}}^{\Lambda})$
and $t \in \mathbb{T}$.
\begin{enumerate}
    \item[(i)] We say that $u$ is \emph{pointed at $t$}
    if there exist $\mathbf{g} \in \Z^J$ and a polynomial
    \[
    F(u_i \mid i \in J_{\mathrm{uf}})
    = \sum_{\mathbf{v} \in \Z_{\geq 0}^{J_{\mathrm{uf}}}}c_{\mathbf{v}}\mathbf{u}^{\mathbf{v}}
    \in \Z[u_i \mid i \in J_{\mathrm{uf}}],
    \qquad c_0 = 1
    \]
    such that
    \[
    u = \mathbf{x}_t^{\mathbf{g}}
    \cdot F(\widehat{y}_{i;t} \mid i \in J_{\mathrm{uf}}).
    \]
    We write $\deg^t u \coloneq \mathbf{g}$
    and call it the \emph{degree} of $u$.
    We write
    \[
    F_u^t(u_i \mid i \in J_{\mathrm{uf}})
    \coloneq F(u_i \mid i \in J_{\mathrm{uf}})
    \]
    and call it the \emph{$F$-polynomial} of $u$.
    \item[(ii)] We say that $u$ is \emph{bipointed at $t$}
    if $u$ is pointed and the $F$-polynomial is of the form
    \[
    F_u^t(u_i \mid i \in J_{\mathrm{uf}})
    = \sum_{0 \leq \mathbf{v} \leq \mathbf{v}_0}c_{\mathbf{v}}\mathbf{u}^{\mathbf{v}},
    \qquad c_0 = c_{\mathbf{v}_0} = 1
    \]
    for some $\mathbf{v}_0 \in \Z_{\geq 0}^{J_{\mathrm{uf}}}$.
    Here, $0 \leq \mathbf{v} \leq \mathbf{v}_0$ means that
    $\mathbf{v}, \mathbf{v}_0 - \mathbf{v} \in \Z_{\geq 0}^{J_{\mathrm{uf}}}$.
    We write
    \[
    \operatorname{codeg}^t u
    \coloneq \deg^t u + \widetilde{B}_t\mathbf{v}_0
    \]
    and call it the \emph{codegree} of $u$.
\end{enumerate}
\end{Def}

\begin{Def}[{\cite{FG09}, \cite{FZ07},
\cite[Definition 3.3]{Cao}}]
\label{Def : tropical transformations and tropical points}
\begin{enumerate}
    \item[(i)] Let $t \in \mathbb{T}$ and $k \in J_{\mathrm{uf}}$.
    We set
    \[
    \phi^Y_{\mu_k t, t} \colon \Z^J \longrightarrow \Z^J;
    \qquad \mathbf{g}^t \longmapsto \mathbf{g}^{\mu_k t},
    \]
    where
    \[
    g_i^{\mu_k t} =
    \begin{cases}
        -g_k^t, & \text{if $i = k$}, \\
        g_i^t + [-b_{ik}^t]_+g_k^t + b_{ik}^t[g_k^t]_+,
        & \text{if $i \neq k$}.
    \end{cases}
    \]
    We set
    \[
    S_Y(\Z^{\max}) \coloneq
    \left\{
    [\mathbf{g}] \coloneq \{\mathbf{g}^t \mid t \in \mathbb{T}\}
    \;\middle|\;
    \begin{gathered}
    \mathbf{g}^{\mu_k t} = \phi^Y_{\mu_k t, t}(\mathbf{g}^t) \\
    \text{for all $t \in \mathbb{T}$, $k \in J_{\mathrm{uf}}$}
    \end{gathered}
    \right\}.
    \]
    \item[(ii)] Let $t \in \mathbb{T}$ and $k \in J_{\mathrm{uf}}$.
    We set
    \[
    \phi^X_{\mu_k t, t} \colon \Z^J \longrightarrow \Z^J;
    \qquad \mathbf{a}^t \longmapsto \mathbf{a}^{\mu_k t},
    \]
    where
    \[
    a_i^{\mu_k t} =
    \begin{cases}
        -a_k^t + \max\left\{
        \displaystyle\sum_{j \in J}[b_{jk}^t]_+a_j^t,
        \displaystyle\sum_{j \in J}[-b_{jk}^t]_+a_j^t
        \right\}, & \text{if $i = k$}, \\
        a_i^t, & \text{if $i \neq k$}.
    \end{cases}
    \]
    We set
    \[
    S_X(\Z^{\max}) \coloneq
    \left\{
    [\mathbf{a}] \coloneq \{\mathbf{a}^t \mid t \in \mathbb{T}\}
    \;\middle|\;
    \begin{gathered}
    \mathbf{a}^{\mu_k t} = \phi^X_{\mu_k t, t}(\mathbf{a}^t) \\
    \text{for all $t \in \mathbb{T}$, $k \in J_{\mathrm{uf}}$}
    \end{gathered}
    \right\}.
    \]
\end{enumerate}
\end{Def}

\begin{Def}[{\cite[Definition 3.4.2]{Qin24}}]
\label{Def : compatibly pointed elements}
Let $u \in \mathcal{U}(\mathcal{S}_{\mathbf{x}}^{\Lambda})$.
We say that $u$ is \emph{compatibly pointed}
if $u$ is pointed at any seed and
\[
\{\deg^t u \mid t \in \mathbb{T}\} \in S_Y(\Z^{\max}).
\]
\end{Def}

It is known that each cluster monomial is compatibly pointed
(see \cite{DWZ10, GHKK18}).

\begin{Lem}[{\cite[Proposition 4.3]{Cao}}]
\label{Lem : tropical points and Lambda}
Let $\mathcal{S}_{\mathbf{x}}^{\Lambda}$
be a $\Lambda$-cluster pattern.
For $[\mathbf{g}] = \{\mathbf{g}^t \mid t \in \mathbb{T}\}
\in S_Y(\Z^{\max})$, we have
\[
\{\Lambda_t\mathbf{g}^t \mid t \in \mathbb{T}\}
\in S_X(\Z^{\max}).
\]
\end{Lem}

\begin{Def}[{\cite{Qin17}}]
\label{Def : injective reachable seeds}
Let
$\mathcal{S}_{\mathbf{x}}^{\Lambda}
= \{(\mathbf{x}_t, \widetilde{B}_t, \Lambda_t)
\mid t \in \mathbb{T}\}$
be a $\Lambda$-cluster pattern.
Let $\operatorname{pr}_{\mathrm{uf}} \colon
\Z^J \longrightarrow \Z^{J_{\mathrm{uf}}}$
be the natural projection.
Then $t \in \mathbb{T}$ is called \emph{injective-reachable}
if there exist $t' \in \mathbb{T}$
and a permutation $\sigma$ of $J_{\mathrm{uf}}$ such that
\[
\operatorname{pr}_{\mathrm{uf}}\deg^t x_{k;t'}
= -f_{\sigma(k)}
\qquad \text{for all $k \in J_{\mathrm{uf}}$},
\]
where $f_{\sigma(k)} \in \Z^{J_{\mathrm{uf}}}$
is the $\sigma(k)$-th unit vector.
\end{Def}

It is known that $t \in \mathbb{T}$ is injective-reachable
if and only if every $t' \in \mathbb{T}$ is injective-reachable
(see \cite{Qin17}).

\begin{Prop}[{\cite[Proposition 3.4.8]{Qin24}}]
\label{Prop : compatible pointedness and bipointedness}
Let $\mathcal{S}_{\mathbf{x}}^{\Lambda}$
be a $\Lambda$-cluster pattern.
Assume that $t \in \mathbb{T}$ is injective-reachable.
If $u \in \mathcal{U}(\mathcal{S}_{\mathbf{x}}^{\Lambda})$
is compatibly pointed, then $u$ is bipointed at any seed.
Moreover, if
$u, u' \in \mathcal{U}(\mathcal{S}_{\mathbf{x}}^{\Lambda})$
are compatibly pointed with the same degree, then
\[
\operatorname{codeg}^{t'}u = \operatorname{codeg}^{t'}u'
\qquad \text{for all $t' \in \mathbb{T}$}.
\]
\end{Prop}


\subsection{Definitions of monoidal categorification of cluster algebras}
\label{ssec : definitions of monoidal categorification}

Let $\mathscr{C}$ be a monoidal category
in Subsection~\ref{ssec : settings in this paper}.
Let $J_{\mathrm{uf}} \subset J$ be finite sets.

\begin{Def}[{\cite[Definition 6.1]{KKOP20}}]
\label{Def : monoidal seed}
A \emph{monoidal seed} in $\mathscr{C}$ is a pair
$\mathscr{S} = (\{M_i\}_{i \in J}, \widetilde{B})$
consisting of a strongly commuting family $\{M_i\}_{i \in J}$
of real modules in $\mathscr{C}$ and an exchange matrix
$\widetilde{B} = (b_{ij})_{(i, j) \in J \times J_{\mathrm{uf}}}$.
\end{Def}

For a monoidal seed $\mathscr{S}$, we set
$\Lambda^{\mathscr{S}} \coloneq (\Lambda(M_i, M_j))_{i, j \in J}$.
By Proposition~\ref{Prop : d vanishing and commutativity},
$\Lambda^{\mathscr{S}}$ is a skew-symmetric integer matrix.

\begin{Def}[{\cite[Definitions 6.3, 6.5]{KKOP20}}]
\label{Def : Lambda admissible monoidal seed}
A monoidal seed
$\mathscr{S} = (\{M_i\}_{i \in J}, \widetilde{B})$
is \emph{$\Lambda$-admissible}
if it satisfies the following conditions:
\begin{enumerate}
    \item[(i)] For any $k \in J_{\mathrm{uf}}$,
    there exists a simple module $M'_k$ in $\mathscr{C}$
    such that there is an exact sequence
    \[
    0 \longrightarrow
    \bigotimes_{i : b_{ik} > 0}M_i^{\otimes b_{ik}}
    \longrightarrow M_k \otimes M'_k
    \longrightarrow
    \bigotimes_{i : b_{ik} < 0}M_i^{\otimes(-b_{ik})}
    \longrightarrow 0,
    \]
    and $M'_k$ commutes with $M_i$ for any $i \neq k$.
    \item[(ii)] $(\widetilde{B}, \Lambda^{\mathscr{S}})$
    is a compatible pair of type $D = 2I_{|J_{\mathrm{uf}}|}$.
\end{enumerate}
\end{Def}

If $\mathscr{S}$ is a $\Lambda$-admissible monoidal seed,
$M'_k$ is uniquely determined by Lemma~\ref{Lem : hd soc injectivity}
For a $\Lambda$-admissible monoidal seed $\mathscr{S}$
and $k \in J_{\mathrm{uf}}$, we set
\[
\mu_k\mathscr{S}
\coloneq
\bigl(\{M_i\}_{i \neq k} \cup \{M'_k\}, \mu_k\widetilde{B}\bigr).
\]

\begin{Prop}[{\cite[Proposition 6.4]{KKOP20}}]
\label{Prop : mutation of Lambda admissible monoidal seeds}
Let $\mathscr{S} = (\{M_i\}_{i \in J}, \widetilde{B})$
be a $\Lambda$-admissible monoidal seed in $\mathscr{C}$
and $k \in J_{\mathrm{uf}}$.
\begin{enumerate}
    \item[(i)] $\mu_k\mathscr{S}$ is a monoidal seed i.e. $M'_k$ is real,
    \item[(ii)] $\Lambda^{\mu_k\mathscr{S}}
    = \mu_k\Lambda^{\mathscr{S}}$.
\end{enumerate}
\end{Prop}

By Proposition~\ref{Prop : mutation of Lambda admissible monoidal seeds}
and the property that compatibility does not change by mutation,
$(\mu_k\widetilde{B}, \Lambda^{\mu_k\mathscr{S}})$
is a compatible pair of type $2I_{|J_{\mathrm{uf}}|}$.

\begin{Def}[{\cite[Definition 7.6]{KKOP24}}]
\label{Def : completely Lambda admissible monoidal seed}
A monoidal seed
$\mathscr{S} = (\{M_i\}_{i \in J}, \widetilde{B})$
is \emph{completely $\Lambda$-admissible}
if $\mathscr{S}$ is $\Lambda$-admissible and,
for any finite sequence $(k_1, \ldots, k_\ell)$ in $J_{\mathrm{uf}}$,
$\mu_{k_\ell}\cdots\mu_{k_1}\mathscr{S}$
is well-defined and a $\Lambda$-admissible monoidal seed.
\end{Def}

\begin{Def}[{\cite[Definition 6.7]{KKOP20}}]
\label{Def : Lambda monoidal categorification}
Let $\overline{\mathcal{A}}(\mathcal{S}_{\mathbf{x}}^{\Lambda})$
be a $\Lambda$-cluster algebra with initial $\Lambda$-seed
$(\mathbf{x}_{t_0}, \widetilde{B}_{t_0}, \Lambda_{t_0})$.
We call $\mathscr{C}$ a \emph{$\Lambda$-monoidal categorification}
of $\overline{\mathcal{A}}$ if it satisfies the following conditions:
\begin{enumerate}
    \item[(1)] There exists a $\Z$-algebra isomorphism
    \[
    \iota \colon K(\mathscr{C})_{\Z}
    \xrightarrow{\sim}
    \overline{\mathcal{A}}(\mathcal{S}_{\mathbf{x}}^{\Lambda}).
    \]
    \item[(2)] There exists a completely $\Lambda$-admissible
    monoidal seed
    $\mathscr{S}_0 \coloneq (\{M_i\}_{i \in J}, \widetilde{B})$
    such that
    \[
    \bigl(\{\iota([M_i])\}_{i \in J},
    \widetilde{B}, \Lambda^{\mathscr{S}_0}\bigr)
    = (\mathbf{x}_{t_0}, \widetilde{B}_{t_0}, \Lambda_{t_0}).
    \]
\end{enumerate}
\end{Def}

By Definition~\ref{Def : Lambda admissible monoidal seed}
and Proposition~\ref{Prop : mutation of Lambda admissible monoidal seeds},
for any finite sequence $(k_1, \ldots, k_m)$ of $J_{\mathrm{uf}}$ and
$\mu_{k_m}\cdots\mu_{k_1}\mathscr{S}_0
= (\{M'_i\}_i, \widetilde{B}')$, we have
\[
\bigl(\{\iota([M'_i])\}_i, \widetilde{B}',
\Lambda^{\mu_{k_m}\cdots\mu_{k_1}\mathscr{S}_0}\bigr)
= (\mathbf{x}_t, \widetilde{B}_t, \Lambda_t),
\]
where $t = \mu_{k_m}\cdots\mu_{k_1}t_0$. 

For each $t \in \mathbb{T}$, the real module corresponding to a cluster variable $x_{i;t}$ is denoted by 
$M_{i;t}$, and the set  $\{M_{i;t} \mid i \in J\}$ is called a \emph{monoidal cluster}. 

\begin{Prop}
\label{Prop : cluster monomials and real prime modules}
Let $\mathscr{C}$ be a $\Lambda$-monoidal categorification of
$\overline{\mathcal{A}}(\mathcal{S}_{\mathbf{x}}^{\Lambda})$
and let
\[
\iota \colon K(\mathscr{C})_{\Z}
\xrightarrow{\sim}
\overline{\mathcal{A}}(\mathcal{S}_{\mathbf{x}}^{\Lambda})
\]
be the corresponding isomorphism.
\begin{enumerate}
    \item[(i)] For any cluster monomial
    $\mathbf{x}_t^{\mathbf{a}} \in \overline{\mathcal{A}}$,
    there exists a real module $M \in \mathscr{C}$
    such that $\iota^{-1}(\mathbf{x}_t^{\mathbf{a}}) = [M]$.
    \item[(ii)] For any cluster variable
    $x_{i;t} \in \overline{\mathcal{A}}$,
    there exists a real prime module $M \in \mathscr{C}$
    such that $\iota^{-1}(x_{i;t}) = [M]$.
\end{enumerate}
\end{Prop}

\begin{proof}
(i) For $t = \mu_{k_m}\cdots\mu_{k_1}t_0$,
$\iota([M_{i;t}]) = x_{i;t}$ for any $i \in J$, where
\[
\mu_{k_m}\cdots\mu_{k_1}\mathscr{S}_0
= (\{M_{i;t}\}_i, \widetilde{B}_t).
\]
Since $\{M_{i;t}\}$ is a real commuting family,
$\bigotimes_{i \in J}M_{i;t}^{\otimes a_i}$
is a real simple module and
\[
\iota^{-1}(\mathbf{x}_t^{\mathbf{a}})
= \left[\bigotimes_{i \in J}M_{i;t}^{\otimes a_i}\right].
\]

(ii) It is enough to show that each $M_{i;t}$ is prime.
By Proposition~\ref{Prop : tensor factors in category},
it is enough to see that $[M_{i;t}]$ in $K(\mathscr{C})_{\Z}$
is a prime element.
By \cite{GLSclu}, each $x_{i;t}$ is an irreducible element
in $\overline{\mathcal{A}}(\mathcal{S}_{\mathbf{x}}^{\Lambda})$.
Since
$K(\mathscr{C})_{\Z}
\simeq \Z\bigl[[S_i] \mid 1 \leq i \leq N\bigr]$,
$K(\mathscr{C})_{\Z}$ is a unique factorization domain.
Now we have that $[M_{i;t}]$ in $K(\mathscr{C})_{\Z}$
is a prime element.
\end{proof}

A simple module corresponding to a cluster monomial is called \emph{reachable}. 

\subsection{$\Lambda$-cluster algebra structure on $K(\mathscr{C})_{\Z}$}
\label{ssec : Lambda cluster algebra structure}

Let $\mathscr{C}$, $N$, $S_k$, $T_k^{t_0}$ $(1 \leq k \leq N)$
be as in Subsection~\ref{ssec : settings in this paper}.
Let $J \coloneq \{1, \ldots, N\}$,
$J_{\mathrm{fr}} \coloneq \{k \in J \mid k^+ > N\}$,
and $J_{\mathrm{uf}} \coloneq J \setminus J_{\mathrm{fr}}$.
Let $Q_{t_0}$ be the quiver with vertex set $J$,
frozen vertex set $J_{\mathrm{fr}}$, and the following arrows:
\begin{itemize}
    \item ordinary arrows:
    $s \xrightarrow{|a_{\iota_s\iota_t}|} t$
    if $1 \leq s < t < s^+ < t^+$ and $s^+ \leq N$,
    \item horizontal arrows:
    $s \longrightarrow s^-$
    if $1 \leq s^- < s \leq N$.
\end{itemize}
We denote by $\widetilde{B}_{t_0}$
the exchange matrix corresponding to $Q_{t_0}$.

\begin{Thm}[{\cite[Theorem 8.1]{KKOP24},
\cite[Theorem 11.2.3]{KKKO18}}]
\label{Thm : Lambda monoidal categorification}
The pair
\[
\mathscr{S}_{t_0}
\coloneq (\{T_k^{t_0}\}_{k \in J}, \widetilde{B}_{t_0})
\]
is a completely $\Lambda$-admissible monoidal seed in $\mathscr{C}$.
For a $\Lambda$-cluster pattern $\mathcal{S}_{\mathbf{x}}^{\Lambda}$
with initial $\Lambda$-seed
$(\mathbf{x}_{t_0}, \widetilde{B}_{t_0},
\Lambda^{\mathscr{S}_{t_0}})$,
$\mathscr{C}$ provides a $\Lambda$-monoidal categorification by
\[
\iota \colon K(\mathscr{C})_{\Z}
\xrightarrow{\sim}
\overline{\mathcal{A}}(\mathcal{S}_{\mathbf{x}}^{\Lambda}),
\qquad
[T_k^{t_0}] \longmapsto x_{k;t_0}.
\]
\end{Thm}

\begin{Rem}
\label{Rem : initial monoidal clusters and i boxes}
In \cite[Theorem 7.20]{KKOP24}, the initial monoidal cluster for quantum affine case is
\[
\left\{
L\left(\prod_{\substack{k \leq s \leq N \\ \iota_s = \iota_k}}
Y_{\pi(\iota_s), q^{p_s}}\right)
\;\middle|\; k \in J
\right\},
\]
not $\{T_k^{t_0}\}_{k \in J}$.
However, we can obtain $\mathscr{S}_{t_0}$
from the monoidal seed in \cite[Theorem 7.20]{KKOP24}
by performing the mutation sequence
proposed in \cite[$\S$13]{GLSKac} and relabeling.
The family $\{T_k^{t_0}\}_{k \in J}$ corresponds to
a family of $\mathbf{i}$-boxes
$\{[k^{\min}, k] \mid k \in J\}$,
and we can calculate the exchange matrix
corresponding to the $\mathbf{i}$-boxes by \cite{KK26}. 
By the same argument, we can check that $t_0 \in \mathbb{T}$ is injective reachable. 
\end{Rem}

\subsection{Laurent families}
\label{ssec : Laurent families}

Let $\mathscr{C}$ be as in
Subsection~\ref{ssec : settings in this paper}, and let
\[
\iota \colon K(\mathscr{C})_{\Z}
\xrightarrow{\sim}
\overline{\mathcal{A}}(\mathcal{S}_{\mathbf{x}}^{\Lambda})
\]
be a $\Lambda$-monoidal categorification as in
Subsection~\ref{ssec : Lambda cluster algebra structure}.

For a real commuting family $\{L_i \mid 1 \leq i \leq r\}$
and $\mathbf{a} \in \Z_{\geq 0}^r$, let
\[
L^{\mathbf{a}} \coloneq \bigotimes_{i=1}^r L_i^{\otimes a_i}.
\]

\begin{Def}[{\cite[Definition 3.2]{Laurentfamily}}]
\label{Def : Laurent family}
A finite set of real modules $\{L_i \mid 1 \leq i \leq r\}$
is called a \emph{Laurent family} in $\mathscr{C}$
if it satisfies the following conditions:
\begin{enumerate}
    \item[(i)] $\mathfrak{d}(L_i, L_j) = 0$
    for $1 \leq i, j \leq r$,
    \item[(ii)] it is independent, i.e.,
    if $L^{\mathbf{a}} \simeq L^{\mathbf{b}}$
    for $\mathbf{a}, \mathbf{b} \in \Z_{\geq 0}^r$,
    then $\mathbf{a} = \mathbf{b}$,
    \item[(iii)] if $X \in \Irr\mathscr{C}$ commutes with all $L_i$,
    then there exists $\mathbf{a} \in \Z_{\geq 0}^r$
    such that $X \simeq L^{\mathbf{a}}$.
\end{enumerate}
\end{Def}

\begin{Prop}[{\cite[Proposition 3.6]{Laurentfamily}}]
\label{Prop : monoidal clusters are Laurent families}
Any monoidal cluster in $\mathscr{C}$ is a Laurent family.
\end{Prop}

\begin{Lem}[{\cite[Lemma 2.14]{Laurentfamily}}]
\label{Lem : algebraic independence of Laurent families}
Let $\mathcal{M} = \{L_i \mid 1 \leq i \leq r\}$
be a Laurent family.
Then $\{[L_i] \mid 1 \leq i \leq r\} \subset K(\mathscr{C})_{\Z}$
is algebraically independent over $\Z$.
\end{Lem}

\begin{Prop}[{\cite[Corollary 3.4]{Laurentfamily}}]
\label{Prop : Laurent expansions with respect to a Laurent family}
Let $\mathcal{M} = \{L_i \mid 1 \leq i \leq r\}$
be a Laurent family.
Then
\[
K(\mathscr{C})_{\Z}
\subset \Z\bigl[[L_i]^{\pm 1} \mid 1 \leq i \leq r\bigr]
\]
as subrings of $\operatorname{Frac}K(\mathscr{C})_{\Z}$.
\end{Prop}

\begin{Prop}
\label{Prop : cardinality of Laurent families}
Any Laurent family has exactly $N$ elements.
\end{Prop}

\begin{proof}
Let $\mathcal{M} = \{L_i \mid 1 \leq i \leq r\}$
be a Laurent family.
Then
\[
\Z\bigl[[L_i] \mid 1 \leq i \leq r\bigr]
\subset K(\mathscr{C})_{\Z}
\subset \Z\bigl[[L_i]^{\pm 1} \mid 1 \leq i \leq r\bigr]
\]
holds by
Proposition~\ref{Prop : Laurent expansions with respect to a Laurent family}.
By Lemma~\ref{Lem : algebraic independence of Laurent families},
the transcendence degrees of
\[
\operatorname{Frac}\Z\bigl[[L_i] \mid 1 \leq i \leq r\bigr],
\qquad
\operatorname{Frac}\Z\bigl[[L_i]^{\pm 1} \mid 1 \leq i \leq r\bigr]
\]
over $\Q$ are both $r$.
However, since $K(\mathscr{C})_{\Z}$ is a polynomial ring
in $N$ variables, the transcendence degree of
$\operatorname{Frac}K(\mathscr{C})_{\Z}$ over $\Q$ is $N$.
Then we have $r = N$.
\end{proof}

Let $\mathcal{M} = \{L_i \mid 1 \leq i \leq N\}$
be a Laurent family.
For $X \in \Irr\mathscr{C}$,
there exists $\mathbf{a} \in \Z_{\geq 0}^N$ such that
\[
[L^{\mathbf{a}} \otimes X]
= \sum_{s=1}^{\ell}c_s[L^{\mathbf{b}_s}]
\qquad
\text{for }c_s \in \Z_{>0},\quad
\mathbf{b}_s \in \Z_{\geq 0}^N
\]
by
Proposition~\ref{Prop : Laurent expansions with respect to a Laurent family}.
Since $L^{\mathbf{a}}$ is real,
$L^{\mathbf{a}} \nabla X$ and $L^{\mathbf{a}} \Delta X$ are simple,
and their multiplicities in $L^{\mathbf{a}} \otimes X$ are both $1$.
We take $\mathbf{h}, \mathbf{s} \in \Z_{\geq 0}^N$ such that
\[
L^{\mathbf{a}} \nabla X \simeq L^{\mathbf{h}},
\qquad
L^{\mathbf{a}} \Delta X \simeq L^{\mathbf{s}}.
\]

\begin{Lem}[{\cite[Lemma 3.1]{Laurentfamily}}]
\label{Lem : well defined degrees for Laurent families}
In this setting, $\mathbf{h} - \mathbf{a}$ and
$\mathbf{s} - \mathbf{a}$ in $\Z^N$ are well-defined.
\end{Lem}

\begin{Def}[{\cite[Definition 3.8]{KK19},
\cite[Definition 3.7]{Laurentfamily}}]
\label{Def : degrees with respect to a Laurent family}
Let $\mathcal{M} = \{L_i \mid 1 \leq i \leq N\}$
be a Laurent family and $X \in \Irr\mathscr{C}$.
We take $\mathbf{h}, \mathbf{s}, \mathbf{a} \in \Z_{\geq 0}^N$
as above.
We write
\[
\deg^{\mathcal{M}}X \coloneq \mathbf{h} - \mathbf{a},
\qquad
\operatorname{codeg}^{\mathcal{M}}X
\coloneq \mathbf{s} - \mathbf{a}.
\]
\end{Def}

For a Laurent family
$\mathcal{M} = \{L_i \mid 1 \leq i \leq N\}$, we set
\[
\begin{aligned}
\Lambda_{\mathcal{M}}
&\coloneq (\Lambda(L_i, L_j))_{1 \leq i, j \leq N}, \\
\Lambda^{\infty}_{\mathcal{M}}
&\coloneq (\Lambda^{\infty}(L_i, L_j))_{1 \leq i, j \leq N}, \\
\widetilde{\Lambda}_{\mathcal{M}}
&\coloneq (\widetilde{\Lambda}(L_i, L_j))_{1 \leq i, j \leq N}.
\end{aligned}
\]

\begin{Prop}[{\cite[Lemma 3.6]{KK19},
\cite[Lemma 5.12]{Laurentfamily}}]
\label{Prop : Lambda inequalities for Laurent expansions}
Let $\mathcal{M} = \{L_i \mid 1 \leq i \leq N\}$
be a Laurent family and $X \in \Irr\mathscr{C}$.
Let
\[
[L^{\mathbf{a}} \otimes X]
= \sum_{s=1}^{\ell}c_s[L^{\mathbf{b}_s}]
\qquad\text{for }\mathbf{b}_s \in \Z_{\geq 0}^N,
\]
and let
\[
L^{\mathbf{a}} \nabla X \simeq L^{\mathbf{h}},
\qquad
L^{\mathbf{a}} \Delta X \simeq L^{\mathbf{s}}.
\]
\begin{enumerate}
    \item[(i)] For any $s$, we have
    \[
    \Lambda^{\infty}_{\mathcal{M}}(\mathbf{h} - \mathbf{b}_s) = 0.
    \]
    \item[(ii)] 
    For $1\leq s \leq \ell$ such that $\mathbf{b}_s \neq \mathbf{h}$, 
    \[
    \Lambda_{\mathcal{M}}(\mathbf{h} - \mathbf{b}_s)\in \Z_{\geq 0}^N\setminus \{0\}. 
    \]
    For $1\leq s \leq \ell$ such that $\mathbf{b}_s \neq \mathbf{s}$, 
    \[
    \Lambda_{\mathcal{M}}(\mathbf{b}_s - \mathbf{s})\in \Z_{\geq 0}^N\setminus \{0\}. 
    \]
\end{enumerate}
\end{Prop}

\begin{Lem}[{\cite[Lemma 3.9]{Laurentfamily}}]
\label{Lem : injectivity of degrees for Laurent families}
Let $\mathcal{M} = \{L_i \mid 1 \leq i \leq N\}$
be a Laurent family.
Then
\[
\deg^{\mathcal{M}},\ \operatorname{codeg}^{\mathcal{M}}
\colon \Irr\mathscr{C} \longrightarrow \Z^N
\]
are injective.
\end{Lem}

\begin{Prop}[{\cite{KK19}, \cite[Proposition 5.24]{CFM}}]
\label{Prop : support of Laurent expansions at monoidal clusters}
Let $\mathcal{M}_t = \{M_{i;t} \mid 1 \leq i \leq N\}$
be a monoidal cluster corresponding to $t \in \mathbb{T}$
and $X \in \Irr\mathscr{C}$.
Let
\[
[M_t^{\mathbf{a}} \otimes X]
= \sum_{s=1}^{\ell}c_s[M_t^{\mathbf{b}_s}]
\qquad\text{for }\mathbf{a}, \mathbf{b}_s \in \Z_{\geq 0}^N,
\]
and let
\[
M_t^{\mathbf{a}} \nabla X \simeq M_t^{\mathbf{h}},
\qquad
M_t^{\mathbf{a}} \Delta X \simeq M_t^{\mathbf{s}}.
\]
Then
\[
\mathbf{b}_s - \mathbf{h},\quad \mathbf{s} - \mathbf{b}_s
\in \widetilde{B}_t\Z_{\geq 0}^{J_{\mathrm{uf}}}.
\]
\end{Prop}

By Proposition~\ref{Prop : support of Laurent expansions at monoidal clusters},
Definition~\ref{Def : pointed and bipointed elements}, and
Lemma~\ref{Lem : Lambda inequalities for heads and socles},
for $t \in \mathbb{T}$,
$\iota([X]) \in \mathcal{U}(\mathcal{S}_{\mathbf{x}}^{\Lambda})$
is bipointed at $t$, and
\[
\deg^{\mathcal{M}_t}X = \deg^t\iota([X]),
\qquad
\operatorname{codeg}^{\mathcal{M}_t}X
= \operatorname{codeg}^t\iota([X]).
\]

\begin{Prop}[{\cite[Theorem 3.16]{KK19}, \cite{Qin17}}]
\label{Prop : compatible pointedness of simple modules}
For $X \in \Irr\mathscr{C}$,
$\iota([X]) \in \mathcal{U}(\mathcal{S}_{\mathbf{x}}^{\Lambda})$
is compatibly pointed.
\end{Prop}

\begin{Prop}
\label{Prop : denominators and commutativity for Laurent families}
Let $\mathcal{M} = \{L_i \mid 1 \leq i \leq N\}$
be a Laurent family and $X \in \Irr\mathscr{C}$.
Let
\[
[X] = \sum_{s=1}^{\ell}c_s\prod_{i=1}^N[L_i]^{\mathbf{b}(s)_i}
\]
be a Laurent expansion, where $c_s \in \Z_{>0}$
and $\mathbf{b}(s) \in \Z^N$.
For $1 \leq k \leq N$, let
\[
\mathbf{d}(X)_k
\coloneq -\min\{\mathbf{b}(s)_k \mid 1 \leq s \leq \ell\}.
\]
Assume that each $[L_i]$ is a prime element in $K(\mathscr{C})_{\Z}$.
Then $\mathbf{d}(X)_k \leq 0$ if and only if $X$ commutes with $L_k$.
\end{Prop}

\begin{proof}
If $\mathbf{d}(X)_k \leq 0$, then we can prove that
$X$ commutes with $L_k$ by the same proof as
\cite[Lemma 4.2]{KK19}.

Assume that $X$ commutes with $L_k$.
By Lemma~\ref{Lem : decrease of d}, there exists $n \in \Z_{\geq 0}$
such that any simple subquotient of
\[
\left(\bigotimes_{i\ne k}L_i^{\otimes n}\right)\otimes X
\]
commutes with $L_i$ for all $i \ne k$.
Since
\[
\mathfrak{d}\left(
\left(\bigotimes_{i\ne k}L_i^{\otimes n}\right)\otimes X, L_k
\right)
= \mathfrak{d}(X, L_k) = 0,
\]
any simple subquotient of
$\left(\bigotimes_{i\ne k}L_i^{\otimes n}\right)\otimes X$
also commutes with $L_k$
by Lemma~\ref{Lem : Lambda tensor products and subquotients}.
By the definition of a Laurent family, any simple subquotient of
$\left(\bigotimes_{i\ne k}L_i^{\otimes n}\right)\otimes X$
is a monomial of $\mathcal{M}$.
It means $\mathbf{d}(X)_k \leq 0$.
\end{proof}

\begin{Prop}\label{Prop : frozen commutativity}
  Let $t \in \mathbb{T}$ and $j \in J_{\mathrm{fr}}$. 
  Then the real prime module $M_{j ; t}$ corresponding to a frozen cluster variable commutes with 
  all simple modules in $\mathscr{C}$. 
\end{Prop}
\begin{proof}
  By \cite{GLSclu}, each $[M_{i ; t}]$ for $i \in J$ is a prime element in $K(\mathscr{C})_{\Z}$. 
  The claim follows by Proposition~\ref{Prop : denominators and commutativity for Laurent families} 
  and the fact that 
  \[
  K(\mathscr{C})_{\Z} \subset 
  \Z[[M_{k ; t}]^{\pm 1}, [M_{i ; t}] \mid k \in J_{\mathrm{uf}}, i \in J_{\mathrm{fr}}]. 
  \]
\end{proof}

\begin{Prop}
\label{Prop : degrees of Laurent families form an integral basis}
Let $\mathcal{M} = \{L_i \mid 1 \leq i \leq N\}$
be a Laurent family and $t \in \mathbb{T}$.
Then
\[
\{\deg^t\iota([L_i]) \mid 1 \leq i \leq N\} \subset \Z^N
\]
is a $\Z$-basis of $\Z^N$.
\end{Prop}

\begin{proof}
We set $\mathbf{g}_i \coloneq \deg^t\iota([L_i])$ and
$\mathcal{M}_t \coloneq \{M_{i;t} \mid 1 \leq i \leq N\}$.
First we prove that $\{\mathbf{g}_i\}$ is linearly independent.
Assume
\[
\sum_{i=1}^N a_i\mathbf{g}_i
= \sum_{i=1}^N b_i\mathbf{g}_i
\qquad\text{for }a_i, b_i \in \Z_{\geq 0},\quad a_i b_i = 0.
\]
Then $\deg^{\mathcal{M}_t}L^{\mathbf{a}}
= \deg^{\mathcal{M}_t}L^{\mathbf{b}}$ holds.
Since $\deg^{\mathcal{M}_t}$ is injective,
$L^{\mathbf{a}} \simeq L^{\mathbf{b}}$.
Since $\mathcal{M}$ is independent, $\mathbf{a} = \mathbf{b}$ holds.

Next we prove that $\{\mathbf{g}_i\}$ generates $\Z^N$.
Let
\[
[L^{\mathbf{a}}\otimes M_{k;t}]
= \sum_{s=1}^{\ell}c_s[L^{\mathbf{b}_s}]
\qquad\text{for }\mathbf{a}, \mathbf{b}_s \in \Z_{\geq 0}^N,
\quad c_s \in \Z_{>0}.
\]
Since each $\iota([L_i])$ for $1 \leq i \leq N$
and $\iota([M_{i;t}])$ is pointed at $t$,
so are $\iota([L^{\mathbf{a}}\otimes M_{k;t}])$
and $\iota([L^{\mathbf{b}_s}])$ for $1 \leq s \leq \ell$.
By the positivity of $c_s$, there exists $1 \leq s_0 \leq \ell$
such that
\[
\deg^t\iota([L^{\mathbf{a}}\otimes M_{k;t}])
= \deg^t\iota([L^{\mathbf{b}_{s_0}}]).
\]
Thus we have
\[
\sum_{i=1}^N a_i\mathbf{g}_i + f_k
= \sum_{i=1}^N(\mathbf{b}_{s_0})_i\mathbf{g}_i,
\]
where $f_k$ is the $k$-th unit vector of $\Z^N$.
This means that $\{\mathbf{g}_i \mid 1 \leq i \leq N\}$
generates $\Z^N$.
\end{proof}

